\documentclass[11pt,a4paper]{article}
\usepackage[T1]{fontenc}
\usepackage[utf8]{inputenc}
\usepackage{newtxtext}
\usepackage{amsmath,amsthm}
\usepackage{newtxmath}
\usepackage[margin=26mm,headheight=15pt]{geometry}
\usepackage{microtype,mathtools,booktabs,array,longtable,enumitem}
\usepackage[dvipsnames]{xcolor}
\usepackage{fancyhdr}
\usepackage[most]{tcolorbox}
\usepackage{hyperref}
\hypersetup{colorlinks=true,linkcolor=MidnightBlue,citecolor=MidnightBlue,urlcolor=MidnightBlue,
 pdftitle={Exact Dynamics of a Three-Dimensional Keller Map},
 pdfauthor={ChatGPT; prepared for Vladimir Reshetnikov},
 pdfsubject={Degree growth, Newton faces, shear spectra, and arithmetic escape}}
\usepackage[nameinlink,noabbrev]{cleveref}
\numberwithin{equation}{section}
\newtheorem{theorem}{Theorem}[section]
\newtheorem{proposition}[theorem]{Proposition}
\newtheorem{lemma}[theorem]{Lemma}
\newtheorem{corollary}[theorem]{Corollary}
\theoremstyle{definition}
\newtheorem{definition}[theorem]{Definition}
\newtheorem{example}[theorem]{Example}
\theoremstyle{remark}
\newtheorem{remark}[theorem]{Remark}
\newcommand{\K}{\mathbb K}
\newcommand{\R}{\mathbb R}
\newcommand{\C}{\mathbb C}
\newcommand{\Q}{\mathbb Q}
\newcommand{\Z}{\mathbb Z}
\newcommand{\N}{\mathbb N}
\newcommand{\A}{\mathbb A}
\newcommand{\one}{\mathbf 1}
\newcommand{\dd}{\operatorname{deg}}
\newcommand{\Newt}{\operatorname{Newt}}
\newcommand{\init}{\operatorname{in}}
\newcommand{\conv}{\operatorname{conv}}
\newcommand{\supp}{\operatorname{supp}}
\newcommand{\diag}{\operatorname{diag}}
\newcommand{\lc}{\operatorname{lc}}
\newcommand{\Spec}{\operatorname{Spec}}
\newcommand{\tr}{\operatorname{tr}}
\newcommand{\Id}{\operatorname{Id}}
\newcommand{\calU}{\mathcal U}
\newcommand{\calL}{\mathcal L}
\newcommand{\calH}{\mathcal H}
\DeclareMathOperator{\charac}{char}
\newcommand{\wtag}{\textup{[write]}}% text added when the report was written into ProveIt
\setlist{itemsep=3pt,topsep=5pt}
\setlength{\parskip}{4pt}
\setlength{\parindent}{1.2em}
\emergencystretch=2em
\pagestyle{fancy}
\fancyhf{}
\fancyhead[L]{\small Exact dynamics of a Keller map}
\fancyhead[R]{\small Research article}
\fancyfoot[C]{\thepage}
\renewcommand{\headrulewidth}{0.3pt}
\newtcolorbox{resultbox}{colback=MidnightBlue!4,colframe=MidnightBlue!55,
 boxrule=0.5pt,arc=1mm,left=3mm,right=3mm,top=2mm,bottom=2mm,breakable}

\begin{document}
\hypersetup{pageanchor=false}
\begin{titlepage}
\thispagestyle{empty}
\vspace*{12mm}
{\sffamily\large PROVEIT \; / \; RESEARCH CONTINUATION\par}
\vspace{15mm}
{\Huge\bfseries Exact Dynamics of a\\[3mm] Three-Dimensional\\[3mm] Keller Map\par}
\vspace{9mm}
{\Large Degree growth, Newton faces, shear spectra,\\[2mm]
and arithmetic escape\par}
\vspace{12mm}
{\large Prepared for Vladimir Reshetnikov\par}
{\normalsize AI-assisted research draft \;\textbullet\; 30 September 2026\par}
\vspace{12mm}
\begin{resultbox}
The first dynamical degree of the repository's explicit map is
\[
\boxed{\lambda_1=3+\sqrt{10}}.
\]
Its iterate degrees satisfy an exact second-order recurrence. Polynomial
source shears give an unbounded, explicitly computed spectrum while
preserving the generic degree three and the Keller condition.
\end{resultbox}
\vfill
\noindent\textbf{Status.} Complete mathematical arguments and exact companion
checks are supplied. The new results have not been peer reviewed or
formalized in Lean or Rocq. The original map, its determinant, its
collisions, and the source-shear construction are credited to prior work.
An exhaustive priority claim is not made.\par
\vspace{5mm}
\noindent\textbf{Repository snapshot.}\quad
\texttt{a866ff9a2cdb9c5f436bb1d82ea5dca59dfee38d}\par
\noindent\small Commit timestamp: 1 October 2026, 00:16 UTC
(30 September in America/Los\_Angeles).\par
\vspace{3mm}
\noindent\small\wtag{} Written into ProveIt's research-report collection
(category \texttt{jacobian-conjecture}) from batch~70, manuscript~06, a
single manuscript; provenance and editorial changes are in
Appendix~\ref{kmd:app:provenance}. None of its new results is formalized.
\end{titlepage}
\hypersetup{pageanchor=true}

\begin{abstract}
We study the iterates of the explicit three-dimensional polynomial map
in ProveIt's Jacobian-conjecture development. We prove that their degrees
are $1,7,43,265,\ldots$, satisfy $d_{n+2}=6d_{n+1}+d_n$, and have first
dynamical degree $3+\sqrt{10}$. This establishes the value suggested in
Liam Giannini's publicly available July 2026 record description; it is
not a claim about the contents of an inaccessible full manuscript or
about global priority. The proof gives considerably more than a limit:
all positive weighted degrees are explicit, the upper Newton geometry
has just two vertices per coordinate, and every initial form on the
intervening wall is a monomial times a power of $xz+3y$.

For every nonzero polynomial $h$ of degree $m$, the source shear
$T_h(x,y,z)=(x,y,z+y^2h(xy))$ produces a map $F_h=F\circ T_h$ with
\[
\lambda_1(F_h)=\frac{2m+7+\sqrt{4m^2+32m+61}}2.
\]
Thus one tame left--right equivalence class contains Keller maps of
fixed generic degree three and unbounded first dynamical degree. We
also determine the different degree laws in characteristic three.
Finally, a monomial-dominance theorem constructs a forward-invariant
complex escape region with a precise logarithmic asymptotic. For
integral coefficients it yields a Zariski-dense set of integral initial
points whose arithmetic degree equals the first dynamical degree.
The package contains thirteen groups of exact checks, a source and
claim ledger, a formalization roadmap, and ten further research
questions.
\end{abstract}
\tableofcontents
\newpage

\section{The question, the baseline, and the contributions}\label{kmd:sec:intro}

\subsection{A concrete dynamical question}
A polynomial map can have three quite different kinds of degree. Its
\emph{ordinary degree} is the largest total degree of a component. Its
\emph{generic degree} is the degree of the induced extension of rational
function fields; in characteristic zero it counts the points of a
generic fiber. Its \emph{first dynamical degree} measures exponential
growth of ordinary degrees under iteration. The present article
separates these quantities and computes the third one exactly for a
specific family whose first two are already understood.

Throughout the characteristic-zero parts, $\K$ is a field of
characteristic zero. Put
\begin{align}
 u&=1+xy,& H&=u^2z+y^2(1+3u),\label{kmd:eq:compact1}\\
 F(x,y,z)&=(P,Q,R)=\bigl(uH,\ y+3xH,\ x(5-3u-x^2z)\bigr).
 \label{kmd:eq:F}
\end{align}
This is precisely the map used in the repository's
\texttt{Algebra/JacobianConjecture} project, not a rescaled or permuted
version. The repository supplies Lean and Rocq proofs of its constant
Jacobian determinant $-2$ and an explicit collision \cite{repo}.
Gao's Theorem 3.3 discusses the same map and its generic degree three
\cite{gao}. We rederive the elementary algebraic facts needed here,
rather than treating every neighboring research report as an axiom.

\noindent\wtag{} \emph{Notation against the rest of ProveIt.} The
polynomial written $H$ in \eqref{kmd:eq:compact1} is the one the
\texttt{Algebra/JacobianConjecture} project and the sibling reports
\cite{weighted,arithmetic} write $h=u^2z+y^2(1+3u)$. In this article
$h$ is reserved for the \emph{shear polynomial} $h(t)$ of
\eqref{kmd:eq:shear-def}, as in the source shears $T_h$ of \cite{weighted};
$H$ never denotes a shear, and $h$ never denotes the component factor.
The inverse parameter $\tau=y+1/x$ of \cref{kmd:prop:generic} is the
variable that \cite{arithmetic} writes $t$, a root of the same cubic
\eqref{kmd:eq:cubic}; it is not the shear invariant $xy$, which the project's research notes
and \cite{weighted} write $t$. The letter $L$ is used locally three
times, each defined where it occurs: the field $\K(P,Q,R)$ in the proof
of \cref{kmd:prop:generic}, the linear form $xz+3y$ of \eqref{kmd:eq:B0}, and a
vector of logarithms in \cref{kmd:sec:escape}.
The map \eqref{kmd:eq:F} is the Lean definition \texttt{counterexample}
(\nolinkurl{Algebra/JacobianConjecture/Lean/JacobianConjecture/Counterexample.lean}),
term for term. Its determinant is the Lean theorem
\texttt{jacobianDet\_counterexample}, over every commutative ring, and
the collision \eqref{kmd:eq:collision} is \texttt{collision}, assembled from
\texttt{collision}$_0$\texttt{\_value} and
\texttt{collision}$_1$\texttt{\_value}; the Rocq counterparts in
\texttt{Coq/Counterexample.v} are \texttt{jacobian\_det\_is\_minus\_two}
and \texttt{integral\_collision\_0\_value},
\texttt{integral\_collision\_1\_value}. These are the only statements of
this article that ProveIt has formalized.

Giannini's Zenodo version 1.1 record, dated 20 July 2026, lists the
first six positive-iterate degrees
\[
7,43,265,1633,10063,62011,
\]
and suggests $3+\sqrt{10}$ as the dynamical degree \cite{giannini}. A finite sequence
alone does not prove that assertion: cancellation might occur later.
The central question is whether one can control the leading terms for
\emph{every} iterate. We answer that question affirmatively and make the
control uniform in positive weights and polynomial source shears.

\subsection{Main conclusions}
Write $F^n=F\circ\cdots\circ F$ ($n$ factors), $F^0=\Id$, and
$d_n=\dd(F^n)$. For a dominant polynomial map $G$, define
\begin{equation}
 \lambda_1(G)=\lim_{n\to\infty}\dd(G^n)^{1/n}.
 \label{kmd:eq:lambda-def}
\end{equation}
The usual existence follows by applying subadditivity to
$\log\dd(G^n)$; our formulas prove existence directly in every case
studied below.

\begin{theorem}[Main degree statements]\label{kmd:thm:main}
Over any field of characteristic different from three, the map
\eqref{kmd:eq:F} has
\begin{equation}
 d_0=1,\qquad d_1=7,\qquad d_{n+2}=6d_{n+1}+d_n,
 \qquad \lambda_1(F)=3+\sqrt{10}.
 \label{kmd:eq:main-base}
\end{equation}
Let $0\ne h\in\K[t]$, $m=\dd h\ge0$, and set
\begin{equation}
 T_h(x,y,z)=(x,y,z+y^2h(xy)),\qquad F_h=F\circ T_h.
 \label{kmd:eq:shear-def}
\end{equation}
For $e_n=\dd(F_h^n)$ one has
\begin{align}
 e_0&=1,&e_1&=2m+8,\nonumber\\
 e_{n+2}&=(2m+7)e_{n+1}+(m+3)e_n,&
 \lambda_1(F_h)&=\frac{2m+7+\sqrt{4m^2+32m+61}}2.
 \label{kmd:eq:main-shear}
\end{align}
In characteristic different from two these maps have determinant $-2$
and separable generic degree three. In characteristic two the degree
formulas still hold, but the Keller condition does not.
\end{theorem}

The proof is distributed across \cref{kmd:sec:algebra,kmd:sec:weighted,kmd:sec:shear}.
No numerical extrapolation enters it. Two further conclusions make the
calculation more than a single recurrence: \cref{kmd:sec:newton} computes the
upper Newton geometry of every iterate, and \cref{kmd:sec:escape} turns the
same matrices into an analytic escape theorem and an arithmetic-degree
statement for a Zariski-dense set of initial points.

\subsection{What is inherited and what is established here}
The distinction matters because the repository contains both
kernel-checked mathematics and unrefereed research reports. The
classification report \cite{weighted} already constructs the shears
\eqref{kmd:eq:shear-def}, describes a much larger normalized weighted
class, and gives its ordinary-degree spectrum. The arithmetic report
\cite{arithmetic} already uses a cubic inverse parameter. We do not
claim those constructions, the original counterexample, or its
generic degree as new.

\begin{center}
\small
\begin{tabular}{@{}p{0.29\textwidth}p{0.62\textwidth}@{}}
\toprule
Object or assertion & Role in this article\\
\midrule
$F$, determinant, collision & Known baseline; rechecked and credited.\\
$T_h$ and weighted classification & Known family; no new classification claimed.\\
$3+\sqrt{10}$ & Previously suggested value; an all-iterate proof is given.\\
Weighted degrees and wall forms & Exact theorems proved here for every iterate.\\
Shear spectrum and characteristic three & Uniform formulas proved here.\\
Escape asymptotics and arithmetic points & Proved here under explicit spectral hypotheses.\\
\bottomrule
\end{tabular}
\end{center}

``Proved here'' describes the contents of this article, not an assertion
that no prior author has obtained the result. The public record
underlying the numerical prediction was accessible; its linked PDF
was not successfully retrieved. Consequently, we neither assign a
conjecture number to that manuscript nor assert that the prediction
remained unproved in all literature. The source ledger in
\cref{kmd:app:sources} records this boundary explicitly.

\section{Algebraic baseline and generic degree}\label{kmd:sec:algebra}

\subsection{Expansion and a localized Jacobian computation}
The full expansion is
\begin{align}
 P={}&z+3xyz+3x^2y^2z+x^3y^3z+4y^2+7xy^3+3x^2y^4,
 \label{kmd:eq:P}\\
 Q={}&y+3xz+6x^2yz+3x^3y^2z+12xy^2+9x^2y^3,
 \label{kmd:eq:Q}\\
 R={}&2x-3x^2y-x^3z.\label{kmd:eq:R}
\end{align}
These formulas will provide a finite support certificate for all later
degree arguments.

\begin{proposition}[Known identities, rederived]\label{kmd:prop:baseline}
As an identity over $\Z[x,y,z]$, $\det J_F=-2$. Moreover,
\begin{equation}
 F(-1,1,5)=F(0,-2,-16)=(0,-2,0).
 \label{kmd:eq:collision}
\end{equation}
Thus over a field of characteristic different from two, $F$ is a
noninjective Keller map.
\end{proposition}
\begin{proof}
The collision is direct substitution. For the determinant, work first
in $\Q(x,y,z)$ and use the intermediate coordinates $(x,u,H)$ from
\eqref{kmd:eq:compact1}. Their Jacobian determinant is $xu^2$. In these
coordinates the target is
\begin{equation}
 \left(uH,\ \frac{u-1}{x}+3xH,
 \frac{x(u+1-x^2H)}{u^2}\right).
 \label{kmd:eq:localized-map}
\end{equation}
Indeed,
\[
 u^2(5-3u)+(u-1)^2(1+3u)=u+1.
\]
Differentiating \eqref{kmd:eq:localized-map} with respect to $(x,u,H)$
gives determinant $-2/(xu^2)$. Multiplication by $xu^2$ proves the
claim in the rational function field, hence in the polynomial ring.
Since both sides have integral coefficients, it is an identity over
$\Z$ and specializes to every characteristic. The two source points
in \eqref{kmd:eq:collision} remain distinct over every field, since their
first coordinates differ by one.
\end{proof}

\subsection{The cubic parameter and what it does not measure}
The next calculation is included to make the comparison between first
and generic degree self-contained. The inverse parameter is the one
used in \cite{arithmetic}; it is not the invariant $xy$ occurring in
the shear.

\begin{proposition}[Generic degree]\label{kmd:prop:generic}
If $\charac\K\ne2$, then $F$ is dominant and
\[
 [\K(x,y,z):\K(P,Q,R)]=3,
\]
with separable extension. Every $F_h$ has the same property, and its
$n$th iterate has function-field degree $3^n$.
\end{proposition}
\begin{proof}
Put $\tau=y+1/x$, and for formal target coordinates $(A,B,C)$ write
\begin{equation}
 g_{A,B,C}(T)=CT^3-2T^2+BT-2A.
 \label{kmd:eq:cubic}
\end{equation}
Expansion using \eqref{kmd:eq:P}--\eqref{kmd:eq:R} proves
\begin{equation}
 g_{P,Q,R}(\tau)=0,\qquad g'_{P,Q,R}(\tau)=2/x.
 \label{kmd:eq:cubic-identities}
\end{equation}
The second identity supplies the rational reconstruction
\begin{equation}
 x=\frac2{g'(\tau)},\qquad y=\tau-\frac1x,
 \qquad z=\frac{2-3xy-C/x}{x^2}.
 \label{kmd:eq:reconstruct}
\end{equation}
Here and below identities involving $1/x$ are identities of rational
function fields, not claims about every special fiber.

Let $L=\K(P,Q,R)$. Equations \eqref{kmd:eq:cubic-identities} and
\eqref{kmd:eq:reconstruct} show $\K(x,y,z)=L(\tau)$ and $[L(\tau):L]\le3$.
Thus $\operatorname{trdeg}_\K L=3$, so $P,Q,R$ are algebraically
independent and $F$ is dominant. The polynomial \eqref{kmd:eq:cubic} is
irreducible in $\K[A,B,C,T]$: it is linear in $A$ with invertible
coefficient $-2$, and its quotient ring is $\K[B,C,T]$. As a polynomial
in $T$ it is primitive over $\K[A,B,C]$, since its $T^2$ coefficient
is the unit $-2$. Gauss's lemma makes it irreducible over
$\K(A,B,C)$. Its derivative does not vanish at $\tau$ by
\eqref{kmd:eq:cubic-identities}; hence the extension is separable of degree
three.

The inverse of $T_h$ is $T_{-h}$. Composition with this polynomial
automorphism preserves generic degree and separability. Finally, the
pullback by a dominant map identifies each successive pair in the
function-field tower with the first pair, so the tower law gives
$3^n$.
\end{proof}

\begin{remark}
Over characteristic zero the top dynamical degree is therefore three.
It must not be substituted for $\lambda_1=3+\sqrt{10}$. Nor does the
number of points of a particular exceptional fiber determine either
of these quantities. This article does not compute the second
dynamical degree.
\end{remark}

\noindent\wtag{} Within ProveIt, the cubic \eqref{kmd:eq:cubic} and the
identities \eqref{kmd:eq:cubic-identities} are the inverse cubic and the
derivative identity $g'(t)=2/x$ of \cite{arithmetic}, where they are
credited in turn to earlier external work; generic degree three of $F$
is Gao's \cite{gao}. \Cref{kmd:prop:generic} is therefore a second route
to a known statement for $F$; what it adds is the passage to the shears
$F_h$ and to the iterates.

\subsection{The exceptional algebra in characteristic two}
Although characteristic two is not a Keller case, it is useful to
separate its algebraic degeneration from the different dynamical
degeneration at three.

\begin{proposition}\label{kmd:prop:char2}
In characteristic two the map $F$ is dominant of purely inseparable
function-field degree two. Its total-degree growth remains
\eqref{kmd:eq:main-base}.
\end{proposition}
\begin{proof}
The degree assertion is proved in \cref{kmd:sec:weighted}; its leading
coefficients remain nonzero in characteristic two. For dominance and
the field extension, write $A=P,B=Q,C=R$ and $\tau=y+1/x$. Reduction
of the polynomial identities gives
\[
 B=C\tau^2,\qquad A=\tau(B+\tau+1/x).
\]
Consequently
\[
 \frac1x=\frac A\tau+B+\tau,\qquad
 y=\tau+\frac1x,\qquad z=\frac{C+x^2y}{x^3}.
\]
It follows that $\K(x,y,z)=\K(A,C,\tau)$. These three generators
are algebraically independent because the left side has transcendence
degree three. The subfield $\K(A,C,C\tau^2)$ is rational of the same
transcendence degree, and $B/C$ is not a square in $\K(A,B,C)$, as
its valuation at the prime $B$ is one. Thus adjoining $\tau$ is a
purely inseparable quadratic extension. The determinant is zero by
\cref{kmd:prop:baseline}; no Keller conclusion is intended.
\end{proof}

\section{Every positive weighted degree and the exact recurrence}
\label{kmd:sec:weighted}

\subsection{A composition lemma that rules out cancellation}
For $w=(a,b,c)\in\R_{>0}^3$ and a nonzero polynomial
$f=\sum_\alpha f_\alpha x^{\alpha_1}y^{\alpha_2}z^{\alpha_3}$, define
\[
 \dd_w f=\max_{f_\alpha\ne0}\alpha\cdot w,
 \qquad
 \init_w f=\sum_{\alpha\cdot w=\dd_w f}f_\alpha X^\alpha.
\]
These definitions allow arbitrary positive real weights, not just
integers. Products of nonzero polynomials satisfy
$\dd_w(fg)=\dd_w f+\dd_w g$: the product of their initial forms is
nonzero in the polynomial ring over a field.

\begin{lemma}[Unique outer leading monomial]\label{kmd:lem:composition}
Let $G=(G_1,G_2,G_3)$ have nonzero polynomial components, with weighted
degree vector $d=(\dd_wG_1,\dd_wG_2,\dd_wG_3)$. Suppose a polynomial
$H_i$ has exactly one monomial $c_iX^{\alpha_i}$ of maximal
$d$-weight. Then
\begin{align*}
 \dd_w H_i(G)&=\alpha_i\cdot d,\\
 \init_w H_i(G)&=c_i\prod_{j=1}^3(\init_wG_j)^{\alpha_{ij}}.
\end{align*}
In particular, there is no requirement that the initial forms of the
$G_j$ themselves be monomials.
\end{lemma}
\begin{proof}
The displayed product is nonzero and has the asserted weight. Every
other outer monomial has strictly smaller $d$-weight, hence after
substitution strictly smaller $w$-degree. It cannot cancel any term
of the displayed product.
\end{proof}

The uniqueness requirement concerns the \emph{outer} polynomial.
This distinction is the reason the argument below also works on the
wall where each initial component of $F$ has two terms.

\subsection{Two matrices and a one-step invariant cone}
Set
\begin{equation}
 M=\begin{pmatrix}3&3&1\\3&2&1\\3&0&1\end{pmatrix},\qquad
 N=\begin{pmatrix}2&4&0\\2&3&0\\2&1&0\end{pmatrix},\qquad
 \one=\begin{pmatrix}1\\1\\1\end{pmatrix}.
 \label{kmd:eq:MN}
\end{equation}
The rows of $M$ give the exponents of
$(x^3y^3z,3x^3y^2z,-x^3z)$; those of $N$ give the exponents of
$(3x^2y^4,9x^2y^3,-3x^2y)$. Assume throughout this section that
$\charac\K\ne3$, so these six coefficients are nonzero.

\begin{lemma}[Support and cone certificate]\label{kmd:lem:cone}
For $w=(a,b,c)>0$ write $\delta(w)=a-b+c$. Then
\begin{equation}
 T(w):=\begin{pmatrix}\dd_wP\\\dd_wQ\\\dd_wR\end{pmatrix}
 =Nw+\max\{\delta(w),0\}\one.
 \label{kmd:eq:tropical}
\end{equation}
If $\delta(w)>0$, the maximal monomial in each component is the one
encoded by $M$; if $\delta(w)<0$, it is the one encoded by $N$.
In all cases $\delta(T(w))>0$.
\end{lemma}
\begin{proof}
Inspection of \eqref{kmd:eq:P}--\eqref{kmd:eq:R} shows that every exponent
in the $i$th component is coordinatewise bounded by either the $i$th
row of $M$ or the $i$th row of $N$. Apart from these two endpoints,
each such bound is strict in at least one coordinate. Positive
weights therefore put all other exponents strictly below the larger
endpoint. The endpoint difference is
\[
 M-N=\one\begin{pmatrix}1&-1&1\end{pmatrix},
\]
which proves \eqref{kmd:eq:tropical}. Finally,
\begin{align*}
 \delta(Mw)&=3a+b+c>0,\\
 \delta(Nw)&=2a+2b>0.
\end{align*}
At $\delta(w)=0$ the two vectors agree, and the conclusion is the same.
\end{proof}

\begin{theorem}[All positive weighted degrees]\label{kmd:thm:weighted}
For every $w\in\R_{>0}^3$ and $n\ge1$,
\begin{equation}
 D_n(w):=\begin{pmatrix}\dd_w(F^n)_1\\\dd_w(F^n)_2\\\dd_w(F^n)_3\end{pmatrix}
 =M^{n-1}\bigl(Nw+\max\{a-b+c,0\}\one\bigr).
 \label{kmd:eq:all-weighted}
\end{equation}
\end{theorem}
\begin{proof}
The first step is \cref{kmd:lem:cone}. Its degree vector belongs to the
positive cone $\delta>0$, where each component of the next outer copy
of $F$ has a unique maximal monomial encoded by $M$. Apply
\cref{kmd:lem:composition} to obtain $D_{n+1}=MD_n$. The same cone lemma
keeps every subsequent vector in $\delta>0$, completing the induction.
\end{proof}

This proof uses only a finite support comparison, two strict linear
inequalities, and an induction. In particular, it proves that no
later cancellation can invalidate the observed degree pattern.

\subsection{Ordinary degrees, generating functions, and a closed form}
For ordinary degree $w=\one$, the initial gap is $1>0$, so
\begin{equation}
 D_n(\one)=M^n\one\qquad(n\ge0).
 \label{kmd:eq:ordinary-matrix}
\end{equation}
The resulting vectors begin as follows.
\begin{center}
\begin{tabular}{@{}rrrr@{}}
\toprule
$n$ & $\dd(F^n)_1$ & $\dd(F^n)_2$ & $\dd(F^n)_3$\\
\midrule
0&1&1&1\\
1&7&6&4\\
2&43&37&25\\
3&265&228&154\\
4&1633&1405&949\\
5&10063&8658&5848\\
6&62011&53353&36037\\
\bottomrule
\end{tabular}
\end{center}
The first coordinate is the largest for $n\ge1$: the differences
between the first and second rows of $M$, and between the second
and third rows, are $(0,1,0)$ and $(0,2,0)$ respectively.

\begin{theorem}[Exact scalar laws]\label{kmd:thm:recurrence}
Let $d_n=\dd F^n$. Then
\begin{equation}
 \sum_{n\ge0}d_nz^n=\frac{1+z}{1-6z-z^2}.
 \label{kmd:eq:gf}
\end{equation}
Writing $s=\sqrt{10}$ and $\lambda_\pm=3\pm s$, one has
\begin{equation}
 d_n=\frac{(s+4)\lambda_+^n+(s-4)\lambda_-^n}{2s}.
 \label{kmd:eq:closed}
\end{equation}
In particular $\lambda_1(F)=\lambda_+$.
\end{theorem}
\begin{proof}
Direct calculation gives
\[
 \det(tI-M)=t(t^2-6t-1),\qquad
 \operatorname{im}M=E:=\{(a,b,c):2a-3b+c=0\}.
\]
The matrix has rank two. On $E$ it satisfies $M^2=6M+I$; this follows
either by restriction or from $M(M^2-6M-I)=0$ and the fact that
$M$ is invertible on its image. Since $\one\in E$, every coordinate
of $M^n\one$ satisfies the recurrence from its initial index $n=0$.
For the first coordinate the initial values are $1,7$, giving
\eqref{kmd:eq:gf} and then \eqref{kmd:eq:closed}. The coefficient of
$\lambda_+^n$ is positive, $\lambda_+>1$, and $|\lambda_-|<1$, so
taking $n$th roots proves the limit.
\end{proof}

For completeness, the other two coordinate generating functions are
\[
 \sum_{n\ge0}\dd(F^n)_2z^n=\frac1{1-6z-z^2},\qquad
 \sum_{n\ge0}\dd(F^n)_3z^n=\frac{1-2z}{1-6z-z^2}.
\]
The small negative conjugate $3-\sqrt{10}$ will reappear in the
transverse contraction governing escape to infinity.

\section{Upper Newton geometry and exact wall forms}\label{kmd:sec:newton}

\subsection{What is computed}
The Newton polytope of a polynomial is the convex hull of all its
exponent vectors. Positive weighted degrees see only its upper
faces. It is useful to encode exactly that information by the
\emph{downward completion}
\begin{equation}
 \mathcal P^\downarrow(f)=\Newt(f)-\R_{\ge0}^3.
 \label{kmd:eq:downward}
\end{equation}
This is not the full Newton polytope. Its support function is finite
on the nonnegative orthant and agrees there with that of
$\Newt(f)$. Calling it an upper Newton object refers to its maximal
faces, even though its completion extends downward.

\begin{theorem}[Two vertices determine the upper geometry]\label{kmd:thm:newton}
Assume $\charac\K\ne3$. For $n\ge1$ and $i\in\{1,2,3\}$ put
\begin{equation}
 v^+_{n,i}=(M^n)_{i,*},\qquad
 v^-_{n,i}=(M^{n-1}N)_{i,*},\qquad k_n=M^{n-1}\one.
 \label{kmd:eq:vertices}
\end{equation}
Then
\begin{equation}
 \mathcal P^\downarrow((F^n)_i)
 =\conv\{v^+_{n,i},v^-_{n,i}\}-\R_{\ge0}^3,
 \label{kmd:eq:newton-theorem}
\end{equation}
and
\begin{equation}
 v^+_{n,i}-v^-_{n,i}=k_{n,i}(1,-1,1).
 \label{kmd:eq:vertex-diff}
\end{equation}
Thus the only wall for strictly positive weights is $a+c=b$,
independently of the iterate.
\end{theorem}
\begin{proof}
By \cref{kmd:thm:weighted}, the support function for positive weights is
\[
 (M^{n-1}Nw)_i+k_{n,i}\max\{a-b+c,0\}
 =\max\{v^+_{n,i}\cdot w,v^-_{n,i}\cdot w\}.
\]
Here $k_{n,i}>0$. Off the wall, successive applications of
\cref{kmd:lem:composition} show that the initial form is a nonzero
monomial of exponent $v^+_{n,i}$ or $v^-_{n,i}$ respectively. Both
endpoints therefore occur in the polynomial. Equality of the
support functions extends to nonnegative weights by continuity.
A closed downward-completed convex polyhedron is determined by these
support values: a separating linear functional with finite upper
bound must have nonnegative coefficients. This proves
\eqref{kmd:eq:newton-theorem}. Equation \eqref{kmd:eq:vertex-diff} follows
by multiplying $M-N=\one(1,-1,1)$ on the left by $M^{n-1}$.
\end{proof}

An equivalent elementary version of the last separation step is the
following. If a support exponent lay outside the right side of
\eqref{kmd:eq:newton-theorem}, a nonnegative weight would strictly
separate it. Perturbing this weight slightly into the positive
orthant would preserve strict separation, contradicting
\cref{kmd:thm:weighted}.

\subsection{The entire initial form on the wall}
Put
\begin{equation}
 L=xz+3y,\qquad
 B_0=\begin{pmatrix}2&3&0\\2&2&0\\2&0&0\end{pmatrix}.
 \label{kmd:eq:B0}
\end{equation}
The notation $B_0$ is an exponent matrix, not the second target
coordinate from \eqref{kmd:eq:cubic}. For a positive wall weight
$a+c=b$, the first initial forms are
\begin{equation}
 \init_w F=(x^2y^3L,\ 3x^2y^2L,\ -x^2L).
 \label{kmd:eq:wall-first}
\end{equation}
They need not remain monomials, but their propagation is nevertheless
exact.

\begin{theorem}[Wall factorization at every iterate]\label{kmd:thm:wall}
Let $C_1=(1,3,-1)$ and recursively define nonzero coefficients
\begin{equation}
 C_{n+1,i}=C_{1,i}\prod_{j=1}^3C_{n,j}^{M_{ij}}.
 \label{kmd:eq:coefficient-rec}
\end{equation}
If $w>0$ and $a+c=b$, then for every $n\ge1$,
\begin{equation}
 \init_w(F^n)_i
 =C_{n,i}\,
 x^{(M^{n-1}B_0)_{i1}}
 y^{(M^{n-1}B_0)_{i2}}
 L^{(M^{n-1}\one)_i}.
 \label{kmd:eq:wall-all}
\end{equation}
\end{theorem}
\begin{proof}
The case $n=1$ is \eqref{kmd:eq:wall-first}. After that first step the
degree vector lies strictly in the plus cone. Therefore the unique
outer initial monomials have exponents $M_{i,*}$ and coefficients
$C_{1,i}$. The composition lemma multiplies the exponent matrix by
$M$, the exponent vector of $L$ by $M$, and the coefficient vector
according to \eqref{kmd:eq:coefficient-rec}. This is exactly the stated
induction. The coefficients remain nonzero because $3\ne0$ and the
coefficient ring is a field.
\end{proof}

For example, $C_2=(-27,-27,1)$. Formula \eqref{kmd:eq:wall-all} describes
all cancellation-sensitive leading terms along this wall. In
positive characteristic some interior binomial coefficients in
$L^k$ can vanish, but neither endpoint of that power vanishes when
$3\ne0$. The upper-geometry theorem is therefore unchanged.

\subsection{A robust support class}
The exact growth is not specific to all sixteen coefficients in
\eqref{kmd:eq:P}--\eqref{kmd:eq:R}. It survives arbitrary changes below the
relevant upper faces.

\begin{corollary}[Coefficient robustness]\label{kmd:cor:robust}
Let $G=(G_1,G_2,G_3)$ be a polynomial tuple over any field. Suppose
\[
 \supp(G_i)\subseteq
 \bigl(\conv\{M_{i,*},N_{i,*}\}-\R_{\ge0}^3\bigr)\cap\N^3
\]
and the coefficient of $X^{M_{i,*}}$ is nonzero for every $i$.
Then the ordinary coordinate-degree vector of $G^n$ is $M^n\one$.
If $G$ is dominant, $\lambda_1(G)=3+\sqrt{10}$.
\end{corollary}
\begin{proof}
For every positive weight in the plus cone, $M_{i,*}$ is the unique
maximizer on the specified support region. The ordinary initial
weight belongs to that cone, and $M$ preserves it. The proof of
\cref{kmd:thm:weighted} therefore applies without change. Dominance is
needed only to use the conventional dynamical-degree terminology,
not for the degree-vector induction.
\end{proof}

In particular, in characteristic different from three the same law
holds after independent invertible affine changes in each source
coordinate and in each target coordinate, with no coordinate mixing
or permutation. Expanding such a substitution moves exponents only
downward from existing ones, and the plus vertex retains a nonzero
coefficient. These source--target modifications need not be
conjugacies. The result is a support statement, not an assertion that
first dynamical degree is invariant under arbitrary left--right
equivalence.

\begin{lemma}[Weighted eigenvector principle]\label{kmd:lem:eigen}
More generally, suppose a polynomial tuple $G$ has unique
$v$-leading monomials $c_iX^{A_{i,*}}$, where $v>0$,
$A\in M_d(\N)$, and $Av=\lambda v$ with $\lambda>1$. Then
\[
 \dd_v(G^n)_i=\lambda^n v_i.
\]
Its ordinary degree has exponential rate $\lambda$.
\end{lemma}
\begin{proof}
The degree vector after each step is a positive scalar multiple of
$v$, so the unique outer leading monomials remain unique. Induction
using \cref{kmd:lem:composition} proves the weighted formula. If
$v_{\min}=\min_i v_i$ and $v_{\max}=\max_i v_i$, then
\[
 v_{\min}\dd f\le\dd_v f\le v_{\max}\dd f
\]
for every nonzero polynomial $f$. These fixed multiplicative
comparison constants disappear on taking $n$th roots.
\end{proof}

This principle is a general leading-term method. Our contribution is
the explicit invariant cones, matrices, and consequences for the
present maps, rather than a claim to have invented weighted degrees.

\section{The complete polynomial-shear spectrum}\label{kmd:sec:shear}

\subsection{Why a shear is dynamically nontrivial}
The transformation $T_h$ in \eqref{kmd:eq:shear-def} is triangular, with
Jacobian one and inverse $T_{-h}$. Thus $F_h=F\circ T_h$ is
left--right equivalent to $F$. It is usually not conjugate to $F$:
iteration produces
\[
 F_h^n=(F\circ T_h)\circ\cdots\circ(F\circ T_h),
\]
not $F^n\circ T_h$ and not $T_h^{-1}\circ F^n\circ T_h$.
This is the distinction that permits unbounded dynamical degrees.
The repository already uses these shears in its weighted
classification \cite{weighted}; here they are treated as dynamical
systems in the fixed displayed coordinates.

Let $0\ne h(t)\in\K[t]$, let $m=\dd h$, and write
$\eta=\lc(h)$. Substitution in \eqref{kmd:eq:F} gives
\begin{align}
 P_h&=P+(1+xy)^3y^2h(xy),\nonumber\\
 Q_h&=Q+3x(1+xy)^2y^2h(xy),\label{kmd:eq:shear-components}\\
 R_h&=R-x^3y^2h(xy).\nonumber
\end{align}
In characteristic different from three the candidate top monomials
are
\begin{equation}
 \eta x^{m+3}y^{m+5},\qquad
 3\eta x^{m+3}y^{m+4},\qquad
 -\eta x^{m+3}y^{m+2}.
 \label{kmd:eq:shear-top}
\end{equation}
Their exponent matrix is
\begin{equation}
 S_m=\begin{pmatrix}
 m+3&m+5&0\\m+3&m+4&0\\m+3&m+2&0
 \end{pmatrix}.
 \label{kmd:eq:S}
\end{equation}

\begin{lemma}[Shear cone]\label{kmd:lem:shear-cone}
For positive weights $w=(a,b,c)$ put
\[
 \Delta_m(w)=ma+(m+2)b-c.
\]
If $\Delta_m(w)>0$, the monomials in \eqref{kmd:eq:shear-top} are the
unique maximal monomials of the corresponding components. Moreover,
\begin{equation}
 \Delta_m(S_mw)=(2m+1)(m+3)a+(2m^2+10m+6)b>0.
 \label{kmd:eq:shear-gap}
\end{equation}
\end{lemma}
\begin{proof}
Every exponent in the added shear part of a component is
coordinatewise bounded by the corresponding row of $S_m$. Only the
highest-degree term of $h$ and the highest power of $xy$ give that
row. The difference between a row of $S_m$ and the corresponding
row of $M$ is $(m,m+2,-1)$; the difference from the row of $N$ is
$(m+1,m+1,0)$. The first difference has positive $w$-weight by
hypothesis, and the second does so because $w>0$. The support
certificate for $F$ now proves uniqueness. Direct multiplication
gives \eqref{kmd:eq:shear-gap}.
\end{proof}

\begin{theorem}[All-iterate shear degree law]\label{kmd:thm:shears}
Assume $\charac\K\ne3$ and $h\ne0$ has degree $m\ge0$. For every
$n\ge0$,
\begin{equation}
 \begin{pmatrix}\dd(F_h^n)_1\\\dd(F_h^n)_2\\\dd(F_h^n)_3\end{pmatrix}
 =S_m^n\one.
 \label{kmd:eq:shear-vector}
\end{equation}
The ordinary degrees satisfy \eqref{kmd:eq:main-shear}, and
\begin{equation}
 \sum_{n\ge0}\dd(F_h^n)z^n
 =\frac{1+z}{1-(2m+7)z-(m+3)z^2}.
 \label{kmd:eq:shear-gf}
\end{equation}
They depend on $h$ only through its degree, not through its other
coefficients or its nonzero leading coefficient.
\end{theorem}
\begin{proof}
The ordinary initial weight satisfies $\Delta_m(\one)=2m+1>0$.
Apply \cref{kmd:lem:shear-cone,kmd:lem:composition} inductively to obtain
\eqref{kmd:eq:shear-vector}. The first coordinate is the largest at
every positive step, since consecutive row differences are
$(0,1,0)$ and $(0,2,0)$.

One has
\[
 \det(tI-S_m)=t\bigl(t^2-(2m+7)t-(m+3)\bigr).
\]
The image is the same plane $E=\{2a-3b+c=0\}$ as before. Its rank
is two, and on $E$ the matrix satisfies
$S_m^2=(2m+7)S_m+(m+3)I$. Since $\one\in E$, this gives the
second-order recurrence from $n=0$. The initial vector is
$(2m+8,2m+7,2m+5)$, yielding \eqref{kmd:eq:shear-gf}.

The two nonzero eigenvalues are
\[
 \lambda_m=\frac{2m+7+\sqrt{4m^2+32m+61}}2,
 \qquad \mu_m=-\frac{m+3}{\lambda_m}.
\]
They satisfy $\lambda_m>m+3$ and $-1<\mu_m<0$. The scalar
recurrence has a positive coefficient of $\lambda_m^n$, because
$e_0=1$, $e_1=2m+8>0$, and $\mu_m<0$. Its exponential rate is
therefore $\lambda_m$.
\end{proof}

\begin{center}
\small
\begin{tabular}{@{}lrrrl@{}}
\toprule
Shear & $\dd F_h$ & $\dd F_h^2$ & $\dd F_h^3$ & $\lambda_1(F_h)$\\
\midrule
$h=0$ &7&43&265&$3+\sqrt{10}$\\
$\dd h=0$ &8&59&437&$(7+\sqrt{61})/2$\\
$\dd h=1$ &10&94&886&$(9+\sqrt{97})/2$\\
$\dd h=2$ &12&137&1567&$(11+\sqrt{141})/2$\\
$\dd h=3$ &14&188&2528&$(13+\sqrt{193})/2$\\
\bottomrule
\end{tabular}
\end{center}

The zero polynomial is a separate case, not a polynomial of degree
zero. In particular, taking a nonzero constant shear coefficient
arbitrarily close to zero over $\C$ does not preserve the first
dynamical degree at the limiting zero coefficient.

\noindent\wtag{} The case $n=1$ of \cref{kmd:thm:shears}, the coordinate
degrees $(2m+8,2m+7,2m+5)$ of $F_h$ (and $(7,6,4)$ for $h=0$), is not
new: it is Corollary~6.4 (``An exact degree spectrum'') of
\cite{weighted}, stated there for the normalized maps $F_{a,b}\circ T_h$.
These differ from $F\circ T_{h'}$ by diagonal scalings of source and
target that replace $h$ by a nonzero multiple of $h(ct)$, $c\ne0$, so
the degrees agree. What is new here is the statement for every
iterate $n\ge2$ and the resulting value of $\lambda_1(F_h)$.

\subsection{Fixed generic degree, unbounded dynamical degree}
\begin{corollary}\label{kmd:cor:unbounded}
Over characteristic zero, one tame left--right equivalence class
contains noninjective Keller maps of separable generic degree three
whose first dynamical degrees are unbounded. More precisely, the
subfamily $\{F_h:h\in\K[t]\}$ realizes exactly
\begin{equation}
 \{3+\sqrt{10}\}\ \cup\
 \left\{\frac{2m+7+\sqrt{4m^2+32m+61}}2:m\in\N\right\}.
 \label{kmd:eq:spectrum}
\end{equation}
\end{corollary}
\begin{proof}
Triangularity of $T_h$ gives $\det J_{F_h}=-2$; generic degree is
three by \cref{kmd:prop:generic}. The collision pulls back to
\begin{align*}
 F_h(-1,1,5-h(-1))
 &=F_h(0,-2,-16-4h(0))\\
 &=(0,-2,0).
\end{align*}
These source points are distinct. The spectrum follows from
\cref{kmd:thm:recurrence,kmd:thm:shears}; its second family is unbounded
because $\lambda_m>2m+7$.
\end{proof}

This does not classify dynamical degrees in the full weighted class,
still less in an entire arbitrary left--right orbit. Diagonal
coordinate scalings preserve the relevant supports, but permutations
and coordinate mixing can change the iteration mechanism. The
precise statement of \cref{kmd:cor:unbounded} is an existence result
inside one equivalence class and an exact spectrum for the specified
source-shear subfamily.

\noindent\wtag{} The Keller property and the collision of $F_h$ in the
proof of \cref{kmd:cor:unbounded} are those of \cite{weighted}, whose
Theorem~6.3 transports collision points by $T_{-h}$; the new content
of the corollary is the unbounded spectrum \eqref{kmd:eq:spectrum}.

\section{A complete degree-growth phase change in characteristic three}
\label{kmd:sec:char3}

All assertions in this section concern iteration in the polynomial
ring over a field of characteristic three. They do not concern the
finite set map obtained by evaluating on a finite field and reducing
polynomial functions modulo $X^q-X$.

\subsection{A simplifying coordinate identity}
Reduction of the components gives
\begin{equation}
 Q=y,\qquad R=-x-x^3z,\qquad P+y^3R=z+y^2.
 \label{kmd:eq:char3-base}
\end{equation}
For the shear family the identities become
\begin{align}
 Q_h&=y,\nonumber\\
 R_h&=-x-x^3z-x^3y^2h(xy),\label{kmd:eq:char3-shear}\\
 P_h+y^3R_h&=z+y^2(1+h(xy)).\nonumber
\end{align}
Thus the second coordinate is fixed under iteration. These exact
identities replace the matrix $M$, whose second leading coefficient
has disappeared.

\begin{theorem}[Characteristic-three phase diagram]\label{kmd:thm:char3}
Let $\charac\K=3$. For $n\ge1$ the coordinate-degree vectors are as
follows.

\smallskip
\noindent\textup{(i)} For $h=0$,
\begin{equation}
 \bigl(\dd(F^n)_1,\dd(F^n)_2,\dd(F^n)_3\bigr)
 =\bigl(7\cdot4^{n-1},1,7\cdot4^{n-1}-3\bigr).
 \label{kmd:eq:char3-zero}
\end{equation}
\noindent\textup{(ii)} For nonzero constant $h$,
\begin{equation}
 \bigl(\dd(F_h^n)_1,\dd(F_h^n)_2,\dd(F_h^n)_3\bigr)
 =\bigl(8\cdot4^{n-1},1,8\cdot4^{n-1}-3\bigr).
 \label{kmd:eq:char3-constant}
\end{equation}
\noindent\textup{(iii)} For $\dd h=m\ge1$,
\begin{equation}
 \bigl(\dd(F_h^n)_1,\dd(F_h^n)_2,\dd(F_h^n)_3\bigr)
 =(a_n,1,a_n-3),
 \label{kmd:eq:char3-vector}
\end{equation}
where
\begin{align}
 a_1&=2m+8,&a_{n+1}&=(m+3)a_n+m+5,\label{kmd:eq:char3-rec}\\
 a_n&=\left(2m+8+\frac{m+5}{m+2}\right)(m+3)^{n-1}
       -\frac{m+5}{m+2}.\label{kmd:eq:char3-closed}
\end{align}
Consequently
\begin{equation}
 \lambda_1(F_h)=
 \begin{cases}
 4,&h=0\text{ or }\dd h=0,\\
 m+3,&\dd h=m\ge1.
 \end{cases}
 \label{kmd:eq:char3-spectrum}
\end{equation}
\end{theorem}
\begin{proof}
For $h=0$, the first degree vector is $(7,1,4)$. For $h\ne0$ of
degree $m\ge0$, it is $(2m+8,1,2m+5)$: the top monomials of
$P_h$ and $R_h$ have coefficients $\eta$ and $-\eta$, which do not
vanish. Suppose at some positive step the degrees are $(a,1,a-3)$.

In the next third coordinate, the term $-x^3z$ has substituted
degree $4a-3$. For a nonzero shear the top term
$-\eta x^{m+3}y^{m+2}$ has substituted degree
$(m+3)a+m+2$. If $m=0$, the difference in favor of $-x^3z$ is
$a-5>0$, since $a\ge8$. If $m\ge1$, the difference in favor of
the shear is
\[
 (m-1)a+m+5>0.
\]
Each winning term is unique. For $h=0$ only the first comparison is
needed, and the linear term $-x$ is smaller.

Use the last identity in \eqref{kmd:eq:char3-shear} after substitution.
The term $-y^3R_h$ has degree three more than the new third
coordinate. The right-side terms have degrees at most $a-3$ and,
when $h\ne0$, $ma+m+2$. They are strictly smaller than that new
first degree. Therefore the new first degree is exactly three more
than the new third degree. This proves preservation of the shape
$(a,1,a-3)$.

The winning term gives $a'=4a$ for $h=0$ and for nonzero constant
$h$, and $a'=(m+3)a+m+5$ for $m\ge1$. The initial values prove
\eqref{kmd:eq:char3-zero}--\eqref{kmd:eq:char3-rec}. Solving the last
linear recurrence over the rational numbers yields
\eqref{kmd:eq:char3-closed}, and taking $n$th roots gives
\eqref{kmd:eq:char3-spectrum}.
\end{proof}

The rational fractions in \eqref{kmd:eq:char3-closed} are expressions
for ordinary integer degree sequences. They are not field elements
in characteristic three; no inversion of $m+2$ in $\K$ is used.
For example, when $m=1$ the formula is $a_n=12\cdot4^{n-1}-2$.

\begin{corollary}[Reduction of an integral shear]
For a fixed $h\in\Z[t]$, reducing modulo three and computing the
actual polynomial $\bar h\in\mathbb F_3[t]$ determines the degree
law completely by \cref{kmd:thm:char3}. Its degree can be smaller than
$\dd h$, and $\bar h$ can be zero. At any prime different from
three where the leading coefficient survives, \cref{kmd:thm:shears}
applies with the original degree.
\end{corollary}

The determinant remains nonzero in characteristic three, and
\cref{kmd:prop:generic} still gives separable generic degree three.
Thus the change from $3+\sqrt{10}$ to $4$ is a genuine change in
iteration complexity, not the loss of dominance or of the Keller
condition.

\section{Escape at infinity and arithmetic degree}\label{kmd:sec:escape}

The degree matrices also control actual orbits whose coordinate
magnitudes lie in a suitable logarithmic cone. A leading monomial
is then not merely largest in a grading: it dominates all other
terms in absolute value, uniformly in their complex arguments.
The small nonleading eigenvalues of the matrices make the
transverse errors contract.

\subsection{A general monomial escape theorem}
For $z=(z_1,\ldots,z_d)\in(\C^*)^d$ set
\[
 \calL(z)=(\log|z_1|,\ldots,\log|z_d|)^\mathsf T.
\]
The following theorem is stated in a dimension-independent form so
that its hypotheses and limitations are visible.

\begin{theorem}[Monomial dominance with contracting transverse spectrum]
\label{kmd:thm:escape}
Let $G=(G_1,\ldots,G_d)$ be a complex polynomial tuple. Suppose
there exist a matrix $A\in M_d(\N)$, a vector $v\in\R_{>0}^d$, a
number $\lambda>1$, and coefficients $c_i\ne0$ such that:
\begin{enumerate}[label=\textup{(\alph*)}]
\item $Av=\lambda v$, and the unique $v$-leading monomial of $G_i$
      is $c_iX^{A_{i,*}}$;
\item $\lambda$ is a simple eigenvalue of $A$, and every other
      eigenvalue has modulus less than one;
\item there is a nonnegative left eigenfunctional $\ell$ with
      $\ell A=\lambda\ell$ and $\ell(v)=1$.
\end{enumerate}
Put $b=(\log|c_1|,\ldots,\log|c_d|)^\mathsf T$ and
\begin{equation}
 L_*=(I-A)^{-1}b.
 \label{kmd:eq:Lstar}
\end{equation}
Then there is a nonempty forward-invariant open set
$\calU\subset(\C^*)^d$ containing all points with logarithms
sufficiently far out in a sufficiently narrow cone around
$L_*+\R_{>0}v$. There is a continuous positive function
$\mathcal G:\calU\to\R_{>0}$ such that
\begin{equation}
 \mathcal G(Gz)=\lambda\mathcal G(z).
 \label{kmd:eq:Gfunctional}
\end{equation}
If $\rho$ is the maximum modulus of the other eigenvalues (zero
when there are none), then for any $q$ with $\rho<q<1$ and every $z\in\calU$,
\begin{equation}
 \calL(G^n z)=\lambda^n\mathcal G(z)v+L_*+O_z(q^n).
 \label{kmd:eq:escape-asymptotic}
\end{equation}
The error can be taken uniform for $z$ in a fixed compact subset
of $\calU$.
\end{theorem}

\begin{proof}
Since no eigenvalue equals one, $L_*$ is defined. Let
$V=\ker\ell$. The restriction $A|_V$ has spectral radius $\rho<1$.
Choose $r_0\in(\rho,1)$ and a norm on $V$ with operator norm at
most $r_0$. This choice is fixed for the construction of $\calU$.
Such a norm exists: for example,
start with any norm and use a convergent weighted sum of
$\|(A|_V)^j e\|$ with ratio chosen strictly between $\rho$ and
$r_0$.

Write logarithmic vectors uniquely as
\begin{equation}
 L=L_*+tv+e,\qquad t=\ell(L-L_*),\quad e\in V.
 \label{kmd:eq:log-decomposition}
\end{equation}
For each nonleading exponent $\alpha$ of $G_i$ there is a strict
gap
\[
 (A_{i,*}-\alpha)\cdot v>0.
\]
There are finitely many such gaps. Hence for sufficiently small
$\varepsilon>0$, whenever $\|e\|<\varepsilon t$ all the
expressions
\[
 (A_{i,*}-\alpha)\cdot(tv+e)
\]
are at least $\kappa_0 t$ for a fixed $\kappa_0>0$. Constants
from $L_*$ and the coefficients can be absorbed into a common
bound. Uniformly in the arguments of the $z_j$, the ratio of each
nonleading term to the leading term is therefore
$O(e^{-\kappa_0t})$.

Choose $R$ so large that these ratios have sum less than $1/2$ in
each component when $t>R$. Then no component of $Gz$ vanishes, and
\begin{equation}
 \calL(Gz)=A\calL(z)+b+E(z),\qquad
 \|E(z)\|\le C e^{-\kappa t}
 \label{kmd:eq:log-error}
\end{equation}
for fixed positive $C,\kappa$. This follows from
$\log|1+r|=O(|r|)$ for $|r|<1/2$; it does not require choosing a
complex logarithm.

Since $AL_*+b=L_*$, the coordinates of the image in
\eqref{kmd:eq:log-decomposition} satisfy
\begin{equation}
 t'=\lambda t+\ell E(z),\qquad
 e'=Ae+(I-v\ell)E(z).
 \label{kmd:eq:te-recurrence}
\end{equation}
Fix $1<\sigma<\lambda$. By increasing $R$, the exponential error
makes $t'\ge\sigma t$ and $\|e'\|<\varepsilon t'$ whenever
$t>R$ and $\|e\|<\varepsilon t$. Indeed, the linear contribution
to $\|e'\|$ is at most $r_0\varepsilon t$, whereas
$\varepsilon t'$ is asymptotic to $\varepsilon\lambda t$;
the difference is a positive linear margin. Thus the set
\begin{equation}
 \calU=\{z\in(\C^*)^d:t(z)>R,\ \|e(z)\|<\varepsilon t(z)\}
 \label{kmd:eq:escape-cone}
\end{equation}
is forward invariant. Shrink $\varepsilon$ as needed to ensure
that all logarithmic coordinates tend to $+\infty$ with $t$;
this is possible because $v>0$.

For $z_n=G^nz$, let $t_n=t(z_n)$ and $E_n=E(z_n)$. Then
$t_n\ge\sigma^nt_0$ and
$\|E_n\|\le C\exp(-\kappa\sigma^nt_0)$. Define
\begin{equation}
 \mathcal G(z)=t_0+\sum_{j=0}^{\infty}
       \lambda^{-j-1}\ell E_j.
 \label{kmd:eq:escape-series}
\end{equation}
The series converges uniformly on compact subsets of $\calU$.
It differs from $t_0$ by an exponentially small quantity as
$t_0\to\infty$, so enlarging $R$ makes it positive. Iterating
\eqref{kmd:eq:te-recurrence} shows that
$\lambda^{-n}t_n\to\mathcal G(z)$ and proves
\eqref{kmd:eq:Gfunctional}. The tail estimate gives
\[
 t_n-\lambda^n\mathcal G(z)
 =O_z\bigl(\exp(-\kappa'\sigma^n)\bigr)
\]
for some $\kappa'>0$.

Finally,
\[
 e_n=A^ne_0+
 \sum_{j=0}^{n-1}(A|_V)^{n-1-j}(I-v\ell)E_j.
\]
For any specified $q\in(\rho,1)$ choose $r\in(\rho,q)$. The
finite-dimensional spectral-radius estimate gives
$\|(A|_V)^n\|\le C_r r^n$ in the fixed norm. Since the forcing
terms decay faster than every geometric sequence, the displayed sum
is $O_z(r^n)$, hence $O_z(q^n)$. This argument does not change
$\calU$. Combining the last two estimates proves \eqref{kmd:eq:escape-asymptotic}, with compact
uniformity following from the same bounds.
\end{proof}

\subsection{Application to the baseline and the shears}
\begin{corollary}[Complex escape for the Keller family]\label{kmd:cor:escape-family}
Over $\C$, both $F$ and every $F_h$ with $h\ne0$ satisfy
\cref{kmd:thm:escape}, with exponent matrix $M$ or $S_m$ respectively
and with the dynamical degrees computed above.
\end{corollary}
\begin{proof}
For $M$ take
\begin{equation}
 v=\begin{pmatrix}
 (2+\sqrt{10})/3\\(7+2\sqrt{10})/9\\1
 \end{pmatrix}.
 \label{kmd:eq:base-eigenvector}
\end{equation}
It is positive and satisfies $Mv=(3+\sqrt{10})v$.
The matrix $M^2$ is strictly positive, so its Perron left
eigenfunctional is positive; alternatively, it can be obtained by
solving the three linear equations directly. The other eigenvalues
are $3-\sqrt{10}$ and zero, both of modulus less than one. The
vector $v$ lies in the plus cone because
$\lambda\delta(v)=\delta(Mv)=3v_1+v_2+v_3>0$.
Thus the required leading monomials are exactly those encoded by
$M$.

For $S_m$, the top-left $2\times2$ block is strictly positive.
Its Perron eigenvector has positive entries, and the last row
extends it to a positive right eigenvector of $S_m$. Its left
Perron eigenfunctional has positive first two entries and zero
third entry. The remaining eigenvalues are $\mu_m\in(-1,0)$
and zero. Equation \eqref{kmd:eq:shear-gap}, evaluated on a positive
eigenvector, shows that $\Delta_m(v)>0$. The unique leading
monomials are consequently \eqref{kmd:eq:shear-top}. All hypotheses
of \cref{kmd:thm:escape} hold.
\end{proof}

For the baseline coefficients $(1,3,-1)$, one obtains the particularly
simple constant correction
\begin{equation}
 L_*=(0,\tfrac12\log3,-\tfrac32\log3)^\mathsf T.
 \label{kmd:eq:base-Lstar}
\end{equation}
Thus, on the escape region, even the subexponential coordinate
ratios are controlled after subtracting the leading ray. The
asymptotic is stronger than merely saying that the orbit is
unbounded.

For $F_h$, set $b_\eta=\log|\eta|$ and $s_3=\log3$. Solving
$(I-S_m)L_*= (b_\eta,b_\eta+s_3,b_\eta)^\mathsf T$ gives
\begin{equation}
 L_*=
 \frac1{3(m+3)}
 \begin{pmatrix}
 -2b_\eta-(m+5)s_3\\
 (m+2)s_3-b_\eta\\
 b_\eta-(4m+11)s_3
 \end{pmatrix}.
 \label{kmd:eq:shear-Lstar}
\end{equation}
Lower coefficients of $h$ influence the escape function and the
size of a valid region, but not the leading exponential rate or
this constant logarithmic correction.

\subsection{Arithmetic degree for many integral points}
For an integral point $p\in\Z^d$, use the logarithmic projective
height
\[
 h([1:p_1:\cdots:p_d])=\log\max\{1,|p_1|,\ldots,|p_d|\}.
\]
For an integral polynomial map $G$ and an integral initial point,
define the arithmetic degree, when the limit exists, by
\begin{equation}
 \alpha_G(p)=\lim_{n\to\infty}
 \max\{1,h([1:G^np])\}^{1/n}.
 \label{kmd:eq:arithmetic-degree}
\end{equation}
It is the \emph{logarithmic} height that appears under the $n$th
root. Taking an $n$th root of the multiplicative height would be
a different and generally divergent expression here.

\begin{theorem}[A Zariski-dense set of maximal arithmetic degree]
\label{kmd:thm:arithmetic}
Let $G=F$, or $G=F_h$ with $0\ne h\in\Z[t]$. There is a
Zariski-dense set of integral initial points $p$ for which
\[
 \alpha_G(p)=\lambda_1(G).
\]
More precisely, every $p\in\calU\cap\Z^3$ in the escape region
from \cref{kmd:cor:escape-family} has this property. With the same
normalization of $v$ and $\mathcal G$ as in \cref{kmd:thm:escape},
\begin{equation}
 h([1:G^np])
 =v_1\mathcal G(p)\lambda_1(G)^n+(L_*)_1+O_p(q^n).
 \label{kmd:eq:height-asymptotic}
\end{equation}
\end{theorem}
\begin{proof}
All iterates remain integral. For the positive right eigenvectors
of both matrices, the first coordinate is strictly largest:
subtracting the first and second eigenvector equations gives
$\lambda(v_1-v_2)=v_2>0$, and subtracting the second and third
gives $\lambda(v_2-v_3)=2v_2>0$. By
\eqref{kmd:eq:escape-asymptotic}, the first logarithmic coordinate
therefore eventually determines the maximum in the height. This
proves \eqref{kmd:eq:height-asymptotic}. Its leading coefficient is
positive, so \eqref{kmd:eq:arithmetic-degree} has limit $\lambda_1(G)$.

It remains to prove that the integral points under consideration
are Zariski dense, rather than merely infinite. For sufficiently
small $\eta_0>0$, every sufficiently large $t$ has the logarithmic
box
\[
 L_*+tv+[-\eta_0t,\eta_0t]^3
\]
contained in the cone defining $\calU$. This follows from the
linear decomposition \eqref{kmd:eq:log-decomposition}: the change in
$t$ and in the transverse part is bounded by a constant times
$\eta_0t$. Choose $\eta_0$ small enough to respect the cone and
to keep $v_i-\eta_0>0$.

Exponentiating this box in the positive real octant gives a
Cartesian box containing arbitrarily many distinct positive
integers in each coordinate interval. A nonzero polynomial of
degree at most $D$ in each variable cannot vanish on a Cartesian
grid with more than $D$ distinct entries in each variable. This
last assertion follows by induction on the number of variables,
using the univariate root bound. Hence no nonzero polynomial can
vanish on all of $\calU\cap\Z^3$, proving Zariski density.
\end{proof}

\begin{remark}[The precise arithmetic boundary]\label{kmd:rem:arithmetic-boundary}
A Zariski-dense \emph{set of initial points} is not the same thing
as a single Zariski-dense orbit. The theorem makes no assertion
that every orbit starting in $\calU$ is Zariski dense. It also does
not assert the equality of arithmetic and dynamical degree for
every rational point with a Zariski-dense orbit. It supplies an
explicitly characterized escape region and proves the equality
for all integral points in that region.
\end{remark}

The construction is nonvacuous at the level of individual points:
rounding each coordinate of
$\exp(L_*+tv)$ to the nearest positive integer places the resulting
point in $\calU$ for all sufficiently large $t$, because the
relative rounding errors tend to zero. The proof by boxes, rather
than this single rounded ray, is what establishes Zariski density.

\section{Verification, scope, and a formalization route}\label{kmd:sec:verification}

\subsection{What the companion computation checks}
The package contains \texttt{code/verify\_results.py}, which was executed
with Python 3.13.5 and SymPy 1.14.0. All thirteen check groups passed;
the recorded output is in \texttt{data/verification.json} and
\texttt{data/verification.txt}. The checks are exact, not floating-point
fits. Their mathematical role is deliberately limited.

\begin{center}
\small
\begin{tabular}{@{}p{0.06\textwidth}p{0.49\textwidth}p{0.34\textwidth}@{}}
\toprule
Group & Object checked & Scope\\
\midrule
1 & Jacobian and integral collision & Exact polynomial identities\\
2 & Cubic equation and derivative identity & Rational-function identities\\
3 & Characteristic-two inverse identities & Exact reduction modulo two\\
4 & Support bounded by the two endpoints & Complete baseline support\\
5 & First wall initial forms & Complete baseline support\\
6 & Matrix polynomial and cone identities & Exact finite identities\\
7 & Ordinary recurrence and degree vectors & First twelve steps\\
8 & Weighted compositions over $\mathbb F_{101}$ & Five weights, three steps\\
9 & Shear characteristic polynomial and gap & Symbolic parameter $m$\\
10 & Shear compositions over $\mathbb F_{101}$ & Four coefficient lists\\
11 & Characteristic-three degree laws & Six cases, four steps\\
12 & Other small characteristics & Primes $2,5,7,11$\\
13 & Spectral and fixed-point identities & Baseline; sample shear checks\\
\bottomrule
\end{tabular}
\end{center}

The composition tests use exact univariate polynomial restrictions
along monomial curves, not full expansion of every multivariate
iterate. Such restrictions supply lower-bound certificates when
they achieve a claimed degree; they can also lose leading terms at
special curves. The selected tests avoid that loss. The upper
bounds and the all-iterate statements come from the proofs, not
from interpolation or from the finite list of tests.

Likewise, the escape theorem is proved analytically. No numerical
orbit simulation is offered as its proof, and no empirical error
bound is substituted for the finite-gap and contraction arguments.
The general shear inequalities are proved for all $m$ in the text;
checking particular eigenvalues in the script does not replace that
argument.

\subsection{A short route to machine-checked degree growth}
The core degree theorem has a relatively small logical footprint.
It need not depend on the repository's classification results or
on the theory of complex dynamics. A useful formalization order is
as follows.

First formalize a support-based weighted degree over multivariate
polynomials over a field. For the ordinary-degree theorem it is
enough to use natural-number weights, finite supports, and the
fact that nonzero products have additive degree. Next prove the
unique-outer-leading-monomial composition lemma. This is the only
place where a cancellation concern needs to be resolved at the
polynomial level.

Then represent $M,N,S_m$ as finite matrices and verify the endpoint
support inequalities from the explicit component polynomials.
The invariance of the plus cone and shear cone is elementary
ordered arithmetic. Induction gives the degree vectors. The
second-order recurrences can be derived directly from
$(M^2-6M-I)\one=0$ and the corresponding identity for $S_m$,
followed by multiplication by matrix powers. This avoids needing
a general spectral theorem in the first formalization pass.

A second pass can establish the generating functions, closed forms,
and limits. A third pass can treat arbitrary real weights and
convex separation for the upper Newton geometry. The escape theorem
is naturally a separate analytic development involving norms,
spectral contraction, and uniformly convergent series.

The inspected baseline Lean source includes \texttt{counterexample}
and \texttt{collision}. Its determinant theorem is
\texttt{jacobianDet\_counterexample}. None of these supplies the new
all-iterate statements in this package. No Lean or Rocq executable
was run for this article, and no uncompiled source is presented as
a formal proof.

\noindent\wtag{} At the placement commit \texttt{51c6943bf} these
declarations are unchanged from the pin: \texttt{counterexample},
\texttt{jacobianDet\_counterexample}, \texttt{collision}$_0$\texttt{\_value},
\texttt{collision}$_1$\texttt{\_value} and \texttt{collision} stand at
lines 95, 134, 151, 158 and 165 of \texttt{Counterexample.lean}, with the
content described above. Nothing in \texttt{Algebra/JacobianConjecture}
concerns iterates, degrees of compositions or dynamics, so the route
above starts from nothing formal beyond the map and its baseline
identities. The placement
of this report beside that Lean and Rocq development confers no formal
status on it.

\subsection{Consequences that are not claimed}
The calculations do not classify unrestricted Keller maps or prove
minimal ordinary degree in dimension three. They do not establish
an algebraically stable compactification, determine the second
dynamical degree, or identify topological entropy with
$\log\lambda_1$. The local escape function is not asserted to be
a global Green function or a global canonical height. The
characteristic-three theorem concerns formal polynomial iteration,
not cycle statistics on finite sets. The upper Newton description
does not determine all lower faces or all coefficients.

Finally, the existence of a Zariski-dense collection of integral
starting points with maximal arithmetic degree does not settle
arithmetic-degree questions for all Zariski-dense rational orbits.
These distinctions are part of the theorem statements, not
post-hoc qualifications to stronger claims.

\section{Further research questions and proposed attacks}\label{kmd:sec:questions}

The following are concrete continuation questions arising from the
proved results. They are proposed research problems, not claims
that each is a previously named open problem in the literature.

\subsection{The second dynamical degree}
Determine $\lambda_2(F)$ and $\lambda_2(F_h)$. We now know
$\lambda_1$ exactly and, in characteristic zero, $\lambda_3=3$.
The matrix $M$ controls coordinate polynomial degrees, but it is
not automatically the action on codimension-two classes. In
particular, taking an exterior power of $M$ without constructing
a suitable geometric model would not prove a formula for
$\lambda_2$. A concrete starting point is to pull back two generic
hyperplanes, compute their intersection degrees after removing
components at infinity, and search for an exact recurrence that
survives saturation.

\subsection{A compactification with stable pullback}
Construct a projective, toric, or explicitly blown-up model on
which the relevant divisor pullback iterates functorially. The
wall $a+c=b$ and the two endpoints in \cref{kmd:thm:newton} provide a
small amount of combinatorial data from which to start. The
challenge is to control centers where the map is indeterminate,
not merely to match the Perron eigenvalue. A successful model
could explain the recurrence geometrically and support a
codimension-two calculation.

\subsection{The global escape locus and the complement of the cone}
Characterize which complex orbits eventually enter $\calU$.
The proof gives a genuine open escape region but does not claim
that every unbounded orbit enters it. The baseline tropical map
has a particularly simple positive-weight chamber structure; the
actual logarithmic dynamics can nevertheless encounter
cancellation near $xz+3y=0$. One approach is to stratify this
cancellation region and derive asymptotic transition laws for its
successive neighborhoods. A useful first target is an explicit
orbit that escapes without eventually approaching the Perron ray,
or a proof that no such orbit exists under a stated nondegeneracy
condition.

\subsection{Arithmetic degree outside the archimedean cone}
For rational initial points whose orbits are Zariski dense, decide
whether the arithmetic degree always exists and equals the value
computed here. An intermediate task is to construct analogous
nonarchimedean escape regions and combine their local height
contributions. Denominators and bad primes must be tracked: the
integral proof avoids them, and therefore does not already give a
number-field theorem. A second useful target is an explicit
classification of exceptional algebraic subvarieties carrying
smaller arithmetic growth.

\subsection{Invariant subvarieties and periodic points}
Determine invariant curves and surfaces, and compare periodic
point counts with degree growth. The rank-two leading matrix and
its image plane describe leading exponents; they do not by
themselves supply a rational first integral or an invariant
hypersurface. Elimination in the ideals generated by $F^n(X)-X$
for small $n$ may reveal patterns, but multiplicities and
positive-dimensional components must be separated from isolated
periodic points. This question would also clarify when points in
the escape region have Zariski-dense individual orbits.

\subsection{Dynamical variation under left--right transformations}
The source shears give an exact unbounded spectrum within one
left--right equivalence class, while the coordinatewise affine
changes of \cref{kmd:cor:robust} preserve the baseline value. Classify
what happens under coordinate permutations, triangular changes
in other directions, and compositions of several shears. A
practical first experiment is to compute the finite list of
leading exponent matrices arising from the six coordinate
permutations, prove invariant cones when they exist, and locate
cases where several tropical chambers recur. This would identify
which changes are dynamically harmless and which genuinely
alter the asymptotic complexity.

\subsection{Beyond the fixed generic-degree-three family}
The public weighted-lift audit \cite{audit} gives ordinary degree
profiles for a different family as its generic degree varies.
Its Proposition 8.3 concerns growth across the family parameter,
not degrees of iterates. Extend the present no-cancellation method
to that family and to the higher-dimensional maps in \cite{gao}.
The target should be a theorem about an explicitly stated support
class or construction, with a verified dominant cone, rather than
an extrapolation from a degree profile. This would compare first
dynamical degree and generic degree across genuinely different
function-field extensions.

\subsection{All specializations and higher-order cancellation walls}
Characteristic three is completely controlled here because the
second coordinate becomes exactly $y$. For more general
coefficient deformations, determine the strata where a leading
coefficient vanishes, where two leading terms cancel after
composition, or where a chamber changes. The wall identity
\eqref{kmd:eq:wall-all} offers a useful test case: its interior support
changes in small characteristic even when its endpoints do not.
A stronger theory would distinguish loss of interior monomials,
loss of a Newton vertex, and change of exponential degree growth.

\subsection{The full Newton polytope and coefficient arithmetic}
Compute the complete Newton polytopes of $F^n$, not just their
downward completions. The wall initial form supplies exact
binomial data on one exposed face, but lower faces may proliferate
and encode phenomena invisible to positive weights. One possible
route is a finite-state recursion for normal cones outside the
positive orthant, accompanied by valuation data for coefficients.
A separate arithmetic question is the growth and divisibility of
the coefficient vector $C_n$ in \eqref{kmd:eq:coefficient-rec}; it is
controlled multiplicatively by the same matrix and is accessible
to exact integer methods.

\subsection{A reusable certified leading-term engine}
Implement a checker that accepts a polynomial tuple, a proposed
leading exponent matrix, and a rational polyhedral cone, and
proves strict dominance plus cone invariance. It should return a
kernel-checked degree recurrence when the matrix admits a small
annihilating polynomial on the initial degree vector. The
baseline and shear family are useful regression tests because
one has a two-chamber first step and the other has an unbounded
symbolic parameter. Extending the checker to nonmonomial initial
forms would make \cref{kmd:thm:wall} a test of more than ordinary
monomial-order reasoning.

\section{Conclusion}
The degree sequence of the repository's three-dimensional map is
not merely compatible with $3+\sqrt{10}$: its entire evolution is
forced by a three-by-three exponent matrix and a strict invariant
cone. The same proof determines every positive weighted degree
and every wall initial form. For polynomial source shears it
gives an exact unbounded spectrum while the generic degree stays
three. In characteristic three a different exact coordinate
identity yields a different, completely explicit phase diagram.

The analytic consequence is equally concrete. The Perron ray is
an attracting direction in logarithmic coordinates for a
nonempty forward-invariant region, with an exponentially small
transverse error after the constant correction. This supplies
maximal arithmetic degree for a Zariski-dense set of integral
starting points. The remaining global dynamical and arithmetic
questions are now separated from the leading-term problem, which
has an exact solution for the displayed family.

\appendix
\section{Source ledger and priority boundary}\label{kmd:app:sources}

\subsection{Repository snapshot}
The ProveIt reference snapshot is
\begin{center}
\texttt{a866ff9a2cdb9c5f436bb1d82ea5dca59dfee38d}.
\end{center}
The main project README, the Jacobian project README, the
weighted-classification README, and the arithmetic-fibers README
were read through the GitHub connector. The original Lean
\texttt{Counterexample.lean} was additionally read at this exact
commit, including the displayed polynomial tuple, the determinant
proof, and the integral-collision declarations. The compiled Lean
and Rocq projects were not rerun.

The classification and arithmetic reports are explicitly marked
unrefereed in their own status descriptions. Our new proofs need
only the displayed definition of $F$ and $T_h$; their correctness
does not depend on the full classification or on the arithmetic
Hasse-principle results in those reports.

\subsection{External primary materials}
Gao's arXiv version 2608.00222v1 was reviewed in its HTML
representation, especially Section 3.4 and Theorem 3.3. Giannini's
Zenodo landing-page description was accessible and supplied the
numerical prediction. Retrieval of its linked PDF failed, and the
linked source repository returned not found through the available
connector. No claim is made about unseen sections of that PDF.

The external audit \cite{audit} was also inspected, in particular
Proposition 8.3. Its formula $(5d-3,5d-4,4)$ concerns component
degrees as the degree $d$ of a construction polynomial varies.
It is not an iterate recurrence. The retrieved file had content
blob identifier
\texttt{a6950580cd582458d1f05ccf259af33e833f9fce}.

Targeted searches for the map together with dynamical-degree and
degree-growth terminology did not establish whether any
unreviewed or later source already proves one of the results
here. Consequently this is an independently derived research
article, not an exhaustive priority determination. In particular,
``the suggested value is proved'' is the supported claim;
``the first proof of a long-standing open conjecture'' is not.

\section{Reproduction and package contents}\label{kmd:app:reproduce}

The article is self-contained LaTeX with an embedded bibliography.
A standard TeX installation providing \texttt{newtxtext},
\texttt{newtxmath}, \texttt{tcolorbox}, and the usual mathematics
packages can build it with two runs of \texttt{pdflatex} (a third
run is harmless). The companion verification uses SymPy 1.14.0:
\begin{verbatim}
python -m pip install -r requirements.txt
python code/verify_results.py
pdflatex -interaction=nonstopmode -halt-on-error article.tex
pdflatex -interaction=nonstopmode -halt-on-error article.tex
\end{verbatim}
A \texttt{Makefile} provides the equivalent targets. Do not run the
Python verifier with optimization flags that disable assertions;
its exact checks intentionally use assertions.

The archive includes the source and compiled PDF, the exact
verification script and its recorded outputs, a compact degree
calculator, and the files \texttt{README.md}, \texttt{SOURCES.md},
\texttt{STATUS.md}, and \texttt{SHA256SUMS}. The source ledger
identifies inherited facts, the status file states the proof
boundaries, and the checksums identify the delivered artifacts.
Font files are not included. Generated TeX auxiliary files and
page-preview images are not part of the delivery.

\noindent\wtag{} The two paragraphs above describe the delivered
package. In ProveIt the layout differs, and some of the commands above
must not be run as printed.
\begin{itemize}
\item In ProveIt the requirements file is
\texttt{data/requirements.txt}, and the Makefile is
\texttt{code/Makefile}. The Makefile's
targets name \texttt{code/verify\_results.py} and \texttt{article.tex}
relative to the package root, so they fail when run from
\texttt{code/}.
\item \texttt{code/verify\_results.py} writes its record, by default, to
\texttt{data/verification.json}: run as printed it overwrites the shipped
record (on Windows with CRLF line endings). Run it on a copy, or pass
\texttt{-{}-output} with a path outside the report. Its recorded console
output \texttt{data/verification.txt} ends with the delivery path
\nolinkurl{/mnt/data/ProveIt_Keller_Dynamics/data/verification.json}.
\item The delivered \texttt{README.md}, the compiled PDF and
\texttt{SHA256SUMS} are not shipped: the README is replaced by the
report's own, \texttt{article.pdf} is a build of the present text, and
the checksum manifest was verified (13 of 13 entries) and retired. The
\texttt{article\_sha256} recorded in \texttt{data/build\_review.json}
is the hash of the delivered PDF, not of \texttt{article.tex}, and that
record describes the delivered 23-page build.
\item \texttt{data/calculator\_checks.json} (930 degree vectors) is
produced by no shipped script; it is kept as delivered.
\end{itemize}

\section{Provenance of this report}\label{kmd:app:provenance}
\wtag{} This report is built from \emph{one} manuscript: manuscript~06
of batch~70 of ProveIt's incoming reports, the archive
\texttt{ProveIt\_Keller\_Dynamics} (inner directory of the same name,
main file \texttt{article.tex}, a 23-page A4 PDF dated 30~September
2026, PDF metadata author ``ChatGPT; prepared for Vladimir
Reshetnikov''), pinned to ProveIt commit
\nolinkurl{a866ff9a2cdb9c5f436bb1d82ea5dca59dfee38d}. The archive
arrived in commit \texttt{1b3960d8a} and was placed in
\texttt{51c6943bf} as a new report of the category
\texttt{jacobian-conjecture}. Between the pin and the placement commit
nothing under \texttt{Algebra/JacobianConjecture} or under the
category's other reports changed. Nothing was merged and nothing was
selected away: every result, proof, example, table, question and
limitation of the manuscript is printed, and the manuscript contributed
everything except the text marked \wtag.

The one choice made at placement was where the manuscript goes. It
studies the map of the formal project and the shears of
\cite{weighted}, but answers no question that \cite{weighted} or
\cite{arithmetic} asks, so it became a report of its own rather than a
Part~III of \cite{weighted}.

The editorial changes are these.
\begin{enumerate}[label=(\arabic*)]
\item Every label carries the prefix \texttt{kmd:}. The manuscript's 87
labels are kept, unchanged after the prefix, and every reference to them
was updated; one label was added, \texttt{kmd:app:provenance}: 88 in
all. The manuscript's macros are kept as delivered.
\item Notes marked \wtag{} were added on the title page, after
\eqref{kmd:eq:F} (notation against the rest of ProveIt and the formalized
declarations), after the remark following \cref{kmd:prop:generic} (prior
sources of the cubic and of generic degree three), after the table
following \cref{kmd:thm:shears} (the case $n=1$ is Corollary~6.4 of
\cite{weighted}), after \cref{kmd:cor:unbounded}, at the end of the
formalization route in \cref{kmd:sec:verification}, and in
Appendix~\ref{kmd:app:reproduce} (shipped layout).
\item The title page no longer creates a PDF page anchor, so the
title page and the first numbered page no longer share one.
\end{enumerate}
No mathematical statement, proof or number of the manuscript was
changed, and the external sources \cite{giannini,gao,audit} were not
re-checked at the write.

\begin{thebibliography}{9}

\bibitem{repo}
Vladimir Reshetnikov and contributors.
\emph{ProveIt: Jacobian-conjecture development}.
Repository documentation and Lean source at commit
\texttt{a866ff9a2cdb9c5f436bb1d82ea5dca59dfee38d}.
\href{https://github.com/VladimirReshetnikov/ProveIt/tree/a866ff9a2cdb9c5f436bb1d82ea5dca59dfee38d/Algebra/JacobianConjecture}{Repository directory at the pinned commit}.

\bibitem{weighted}
ProveIt research-report collection.
\emph{All-Degree Rigidity of a Weighted Keller Class}, Parts I--II.
AI-assisted, unrefereed reports dated 24 and 29 September 2026;
README and status ledger reviewed.
\href{https://github.com/VladimirReshetnikov/ProveIt/tree/a866ff9a2cdb9c5f436bb1d82ea5dca59dfee38d/SetTheory/Cardinals/docs/reports/jacobian-conjecture/weighted-keller-rigidity}{Repository directory at the pinned commit}.

\bibitem{arithmetic}
ProveIt research-report collection.
\emph{Arithmetic local--global fibers}, Parts I--III.
AI-assisted, unrefereed reports; README and notation ledger reviewed.
\href{https://github.com/VladimirReshetnikov/ProveIt/tree/a866ff9a2cdb9c5f436bb1d82ea5dca59dfee38d/SetTheory/Cardinals/docs/reports/jacobian-conjecture/arithmetic-local-global-fibers}{Repository directory at the pinned commit}.

\bibitem{giannini}
Liam Giannini.
\emph{The Alp\"oge--Fable counterexample to the Jacobian Conjecture:
verification, geometry, dynamics, and an equivariant program for
the plane case}.
Zenodo, version 1.1, 20 July 2026.
DOI: \href{https://doi.org/10.5281/zenodo.21461572}{10.5281/zenodo.21461572}.
Public record description reviewed; full PDF not successfully retrieved.

\bibitem{gao}
Shuhong Gao.
\emph{Counterexamples to the Jacobian conjecture in dimensions
greater than two}.
arXiv:2608.00222v1, 31 July 2026, Section 3.4 and Theorem 3.3.
\url{https://arxiv.org/html/2608.00222v1}.

\bibitem{audit}
\texttt{shadybrook/jacobian-counterexample-audit} contributors.
\emph{An Exact Audit of an Announced Three-Dimensional Keller Map}.
AI-assisted working note, stated revision 21 July 2026;
Proposition 8.3 and adjacent discussion reviewed.
\href{https://github.com/shadybrook/jacobian-counterexample-audit/blob/main/paper/main.md}{Public working-note source}.

\end{thebibliography}
\end{document}
